\chapter{Proyectos Aplicados: Elementos de Máquina}

\section{Proyecto Final: Aplicación en Elementos de Máquina}

La finalidad última de esta guía es aplicar todos los conocimientos
adquiridos a proyectos reales de ingeniería mecánica. El diseño
asistido por computadora (CAD) no es un fin en sí mismo, sino una
herramienta fundamental para modelar componentes que deben interactuar
física y mecánicamente entre sí.

\subsection{Normalización de Componentes}

En la industria, rara vez se diseña un tornillo, una chaveta o un
rodamiento desde cero. Estos componentes están \textbf{normalizados}
(IRAM, DIN, ISO), lo que significa que tienen medidas, materiales y
resistencias estándares a nivel mundial.

El desafío del proyectista no es diseñar el componente en sí, sino
saber \textbf{seleccionar} el componente adecuado para la aplicación y
\textbf{representarlo} correctamente en el plano técnico.

\subsection{El Conjunto Mecánico}

Un proyecto CAD de conjunto mecánico implica:
\begin{enumerate}
	\item \textbf{Modelado de piezas individuales:} Crear cada pieza
	      respetando sus tolerancias y ajustes.
	\item \textbf{Ensamblaje virtual:} Unir las piezas en el espacio
	      3D (si el software lo permite) o representar su montaje en
	      vistas 2D (plano de conjunto).
	\item \textbf{Lista de piezas (BOM - Bill of Materials):} Es una
	      tabla indispensable que acompaña al plano de conjunto,
	      enumerando cada pieza, su nombre, cantidad, material y
	      número de norma.
\end{enumerate}

La \textbf{tolerancia de ajuste} (eje-agujero) es un concepto vital:
determina si dos piezas ensambladas tendrán juego (movimiento libre),
estarán ajustadas (fricción) o requerirán presión (interferencia).

La realización de este proyecto final integrará los conocimientos de
vistas, cortes, acotado, capas y bloques para entregar un conjunto de
planos profesionales, listos para la fabricación en el taller.

\section{Cuestionario}
\begin{enumerate}
	\item ¿Por qué es fundamental utilizar elementos de máquina
	      normalizados en lugar de diseñarlos desde cero?
	\item ¿Cómo se representa convencionalmente un tornillo y su rosca
	      en un plano técnico?
	\item ¿Por qué es crítica la tolerancia en el ajuste de un eje y un
	      rodamiento?
	\item ¿Cómo se define el tipo de ajuste (juego, transición,
	      interferencia) entre dos piezas?
	\item ¿Para qué se utiliza una chaveta en un conjunto eje-polea?
	\item Define qué es un elemento de máquina estático vs dinámico.
	\item ¿Cómo se representan las piezas ensambladas en un plano de
	      conjunto?
	\item ¿Qué información técnica mínima debe tener una tabla de lista
	      de piezas (BOM)?
\end{enumerate}

\section{Parte Práctica}
\begin{enumerate}
	\item Dibujar un perno y su tuerca según norma, utilizando bloques
	      para facilitar la inserción.
	\item Dibujar un eje con su chavetero, aplicando las tolerancias
	      necesarias.
	\item Representar un rodamiento simple mediante cortes y secciones.
	\item Realizar el conjunto de un eje, una chaveta y una polea en
	      una vista de corte.
	\item Crear una tabla de lista de piezas (BOM) vinculada a los
	      elementos del conjunto anterior.
	\item Dibujar una pieza mecánica simple (soporte) con tolerancias
	      dimensionales ajustadas.
	\item Realizar las vistas de un cuerpo de válvula con corte total.
	\item Dibujar un engranaje recto básico, incluyendo datos de
	      módulo y número de dientes.
	\item Realizar un plano de montaje completo de un conjunto
	      mecánico (ej: caja reductora simple).
	\item Acotar correctamente todas las piezas del conjunto,
	      asegurando que las medidas coincidan entre piezas
	      ensambladas.
\end{enumerate}
